\section{Graphons, entropy, and large deviations}\label{sec:graphons-entropy}

We use the cut topology for existence and conditional concentration,
and norms of measurable kernels for the perturbation and limit
arguments. We record the relevant continuity properties and entropy
identities here.

\subsection{Graphons and ordered blocks}\label{sec:graphons}

All graphs are finite and simple.  Write $N:=|V(H)|$ and $m:=|E(H)|\ge1$,
and discard isolated vertices of $H$, since they do not change $t(H,W)$.
For $v\in V(H)$,
we write $N_H(v)$ for the set of neighbors of $v$ and $d_v:=|N_H(v)|$ for its
degree, and $\Delta(H)$ for the maximum degree.
Thus $\Delta(H)=1$ exactly when $H$ is a perfect matching.
For a graph $G$ with vertex set $\{1,\ldots,n\}$,
its associated graphon is defined by
\[
  W_G(x,y):=
  \begin{cases}
    1,&\text{if }(\lceil nx\rceil,\lceil ny\rceil)\in E(G),\\
    0,&\text{otherwise}.
  \end{cases}
\]
This representation preserves homomorphism densities:
$t(H,W_G)=t(H,G)$.

We denote Lebesgue measure on $[0,1]$ and its product measures by $\nu$.
Throughout the paper, equalities and inequalities between measurable
functions hold almost everywhere.  The real Banach space
$L^\infty_{\mathrm{sym}}([0,1]^2)$ consists of the essentially bounded
symmetric measurable kernels, identified up to equality almost everywhere.
The graphon space defined in \Cref{sec:main-results} is its subset
\[
  \W_0=\bigl\{W\in L^\infty_{\mathrm{sym}}([0,1]^2):0\le W\le1\bigr\}.
\]
For $F\in L^\infty_{\mathrm{sym}}([0,1]^2)$ and graphons $U,W$, define
\[
  \|F\|_\cut:=\sup_{S,T\subseteq[0,1]}
       \left|\int_{S\times T}F(x,y)\dd x\dd y\right|,
  \qquad
  \delta_\cut(U,W):=\inf_\phi\|U-W^\phi\|_\cut,
\]
where $S,T$ are measurable and $\phi:[0,1]\to[0,1]$ ranges over
measure-preserving bijections modulo null sets, and $W^\phi(x,y):=W(\phi x,\phi y)$.
We call such maps \emph{relabelings}.
The cut norm compares edge densities between measurable vertex sets,
and the infimum removes the dependence on their labels. We call graphons
equivalent when their cut distance is zero; the quotient is the metric
space $\widetilde\W_0$.
This identification is coarser than equality almost everywhere: relabeling
can change a graphon as a measurable function without changing its class in
$\widetilde\W_0$.  For a constant graphon, however,
\[
  \delta_\cut(W,p)=0\quad\Longleftrightarrow\quad W=p
  \quad\text{almost everywhere}.
\]
The densities $e(W)$ and $t(H,W)$ and the entropy cost $I_p(W)$
depend only on the cut-distance equivalence class, so we use the same
notation for them on $\widetilde\W_0$. Kernel norms and perturbations
are taken with vertex labels fixed.
We use the following standard facts about the cut
topology~\cite{LovaszGraphLimits,ChatterjeeVaradhanLDP}.

\begin{lemma}[Compactness and continuity]\label{lem:graphon-compactness}
The space $\widetilde\W_0$ is compact.  For every finite simple graph $H$
and every $p\in(0,1)$, the functions $W\mapsto t(H,W)$ and $W\mapsto e(W)$
are continuous on $\widetilde\W_0$, and $W\mapsto I_p(W)$ is lower
semicontinuous.  In particular, $I_p$ attains its minimum on every
nonempty closed subset of $\widetilde\W_0$.
\end{lemma}

For a graphon $W$, its \emph{row} at $x$ and its
\emph{degree function} are
\[
  W_x(y):=W(x,y),\qquad
  d_W(x):=\int_0^1W_x(y)\dd y.
\]
By Fubini's theorem, $W_x$ is measurable for almost every $x$, and $d_W$
is measurable.  Two rows are regarded as equal if they agree almost
everywhere.

For a measurable row $u:[0,1]\to[0,1]$ and $p\in(0,1)$, we also write
\[
  I_p(u):=\int_0^1J_p\bigl(u(y)\bigr)\dd y
\]
for its entropy cost. With this notation, Fubini's theorem gives
$I_p(W)=\int_0^1I_p(W_x)\dd x$.

For a multipodal graphon, the block structure gives $W_x=b_{ij}$ on $A_j$
for almost every $x\in A_i$.
The blocks, also called \emph{podes}, are ordered when
\begin{equation}\label{eq:ordered-blocks}
  b_{ij}\le b_{i'j}\qquad\text{whenever }i\le i',\ \text{for every }j;
\end{equation}
by symmetry, $(b_{ij})$ is then nondecreasing in both indices.
Equivalently, after discarding a null set of vertex labels, the distinct
row types can be listed as $g_1\le\cdots\le g_M$.
After merging identical rows, their means are strictly increasing, so
these means distinguish the blocks.

\begin{defn}[Ordered structure]\label{def:ordered-structure}
A graphon $W$ has an \emph{ordered structure} if there is a
measure-preserving relabeling $\phi:[0,1]\to[0,1]$ such that, after
modification on a null set, $W^\phi$ is nondecreasing in both variables:
\[
  W^\phi(x,y)\le W^\phi(x',y')
  \qquad\text{whenever }x\le x'\text{ and }y\le y'.
\]
\end{defn}

For a multipodal graphon, this is equivalent to having ordered blocks:
placing the blocks on consecutive intervals makes the graphon
nondecreasing in both variables. In general, the definition allows infinitely
many row types.

\subsection{Entropy and continuity}\label{sec:relative-entropy}

The Bernoulli relative entropy $J_p$ is nonnegative and vanishes only at
$p$.  Thus $I_p(W)\ge0$, with equality exactly when $W=p$ almost everywhere.
We write
\[
  \lo(z):=\log\frac z{1-z},
  \qquad
  \sigma(\xi):=\frac1{1+e^{-\xi}}
\]
for the log-odds function $\lo:(0,1)\to\RR$ and its inverse
$\sigma:\RR\to(0,1)$.  For $0<z<1$,
\begin{equation}\label{eq:entropy-derivatives}
  J_p'(z)=\lo(z)-\lo(p),
  \qquad
  J_p''(z)=\frac1{z(1-z)}\ge4.
\end{equation}
Thus $J_p$ is strictly convex, decreases on $[0,p]$, and increases on
$[p,1]$.  Although $J_p$ is continuous on $[0,1]$, its derivative satisfies
\begin{equation}\label{eq:entropy-endpoint-limits}
  \lim_{z\downarrow0}J_p'(z)=-\infty,
  \qquad \lim_{z\uparrow1}J_p'(z)=+\infty.
\end{equation}
Moving values away from either endpoint can therefore give an
arbitrarily large first-order entropy saving per unit change.
This is the basis of the interiority arguments below.
The inverse log-odds function satisfies
\begin{equation}\label{eq:logistic-derivative}
  0<\sigma'(\xi)=\sigma(\xi)\bigl(1-\sigma(\xi)\bigr)\le\frac14
  \qquad(\xi\in\RR),
\end{equation}
so composition with $\sigma$ is Lipschitz in each $L^q([0,1])$ norm,
$1\le q\le\infty$.

Changing the reference probability changes $J_p$ by an affine function:
for $p,q\in(0,1)$ and $z\in[0,1]$,
\begin{equation}\label{eq:entropy-reference-change}
  J_p(z)-J_q(z)
  =\bigl(\lo(q)-\lo(p)\bigr)z+\log\frac{1-q}{1-p}.
\end{equation}
Consequently, at fixed edge density, the choice of reference probability
does not change the minimizers. The definition of $S_0$ in
\Cref{sec:intro-fixed} also gives
\[
  I_p(W)=-2s(W)-\log(1-p)-e(W)\lo(p).
\]
Taking $p=\eps$ and $e(W)=\eps$ yields
\begin{equation}\label{eq:entropy-identity}
  I_\eps(W)=2\bigl(S_0(\eps)-s(W)\bigr)\qquad\text{whenever }e(W)=\eps.
\end{equation}
Thus, among graphons with edge density $\eps$, minimizing $I_\eps$
is equivalent to maximizing $s$.

We will also compare graphons in $L^1([0,1]^2)$, keeping their vertex
labels fixed.  For graphons $U,W$ and a graph $H$ with $m$ edges,
\begin{equation}\label{eq:density-l1-continuity}
  |t(H,U)-t(H,W)|\le m\|U-W\|_1.
\end{equation}
To see this, order the edges of $H$ and apply the product identity
\begin{equation}\label{eq:product-difference}
  \prod_{i=1}^m a_i-\prod_{i=1}^m b_i
  =\sum_{i=1}^m(a_i-b_i)
       \prod_{k<i}a_k\prod_{k>i}b_k.
\end{equation}
In the density integrals, take $a_i=U(x_v,x_w)$ and $b_i=W(x_v,x_w)$
for the $i$th edge $vw$.  All other factors lie in $[0,1]$, so each
summand contributes at most $\|U-W\|_1$ in absolute value.

On $\W_0$ with the $L^1([0,1]^2)$ topology, relative entropy is continuous:
\begin{equation}\label{eq:entropy-l1-continuity}
  \|W_n-W\|_1\longrightarrow0
  \quad\Longrightarrow\quad I_p(W_n)\longrightarrow I_p(W).
\end{equation}
For every $\eta>0$, uniform continuity of $J_p$ on $[0,1]$ gives
$\delta>0$ such that $|J_p(a)-J_p(b)|\le\eta$ whenever $|a-b|\le\delta$.
Splitting the integral according to this inequality yields
\[
  |I_p(W_n)-I_p(W)|
  \le\eta+\frac{2\max_{z\in[0,1]}J_p(z)}{\delta}\|W_n-W\|_1.
\]
First letting $n\to\infty$ and then $\eta\downarrow0$
proves~\eqref{eq:entropy-l1-continuity}.

\subsection{Large deviations and conditioning}\label{sec:large-deviations}

For fixed $p\in(0,1)$, the Chatterjee--Varadhan
principle~\cite{ChatterjeeVaradhanLDP} is a large-deviation principle for
$W_{G(n,p)}$ in $\widetilde\W_0$, with speed $\binom n2$ and rate function
$I_p$.  For every open set $\mathcal O$ and closed set $\mathcal C$ in
$\widetilde\W_0$, it gives
\begin{align*}
  \liminf_{n\to\infty}\frac1{\binom n2}
    \log\mathbb P(W_{G(n,p)}\in\mathcal O)
    &\ge-\inf_{W\in\mathcal O}I_p(W),\\
  \limsup_{n\to\infty}\frac1{\binom n2}
    \log\mathbb P(W_{G(n,p)}\in\mathcal C)
    &\le-\inf_{W\in\mathcal C}I_p(W).
\end{align*}
For the joint upper-tail constraints, these bounds give the
rate~\eqref{eq:joint-upper-tail-rate}.

Conditional concentration is expressed in terms of the distance to the
optimizer set: for nonempty $\mathcal K\subseteq\widetilde\W_0$, write
\[
  \delta_\cut(U,\mathcal K):=\inf_{W\in\mathcal K}\delta_\cut(U,W).
\]

Let $\mathcal K_{\mathbf H}(p,\mathbf r)$ be the set of minimizers
of~\eqref{eq:joint-variational-problem}, viewed in $\widetilde\W_0$.
For every $\eta>0$,
\begin{equation*}
  \mathbb P\!\left(
    \delta_\cut\bigl(W_{G(n,p)},\mathcal K_{\mathbf H}(p,\mathbf r)\bigr)\ge\eta
    \,\middle|\,
    t(H_j,G(n,p))\ge r_j^{m_j},\, 1\le j\le s
  \right)\rightarrow0
  \qquad (n\to\infty).
\end{equation*}
Both this concentration statement and~\eqref{eq:joint-upper-tail-rate}
follow because the open and closed joint constraint sets have the same
infimum:
replacing a feasible $W$ by $(1-\theta)W+\theta$, with $0<\theta<1$, makes
all inequalities strict, while its cost tends to $I_p(W)$ as $\theta\downarrow0$.
By \Cref{lem:graphon-compactness}, the minimum over feasible graphons
outside any fixed neighborhood of the minimizer set is strictly larger.
Applying the large-deviation upper bound to this restricted set and the
lower bound to the conditioning event proves concentration; see
also~\cite[Theorem~3.1 and its proof]{ChatterjeeVaradhanLDP}.

For fixed densities, let $\mathcal K_H(\eps,\tau)$ be the set of maximizers
of~\eqref{eq:entropy-maximization} in $\widetilde\W_0$.
Choose $G_{n,\delta}$ uniformly among labeled graphs on $n$ vertices satisfying
\[
  |e(W_G)-\eps|<\delta,
  \qquad |t(H,G)-\tau|<\delta.
\]
Then, for every $\eta>0$,
\begin{equation*}
  \lim_{\delta\downarrow0}\limsup_{n\to\infty}
  \mathbb P\!\left(
    \delta_\cut\bigl(W_{G_{n,\delta}},\mathcal K_H(\eps,\tau)\bigr)\ge\eta
  \right)=0.
\end{equation*}
The uniform law here is that of $G(n,1/2)$ conditioned on the density
window.  Since $I_{1/2}(W)=\log2-2s(W)$, its rate minimizers on the exact
constraint set are precisely $\mathcal K_H(\eps,\tau)$, and the same
compactness and conditioning argument applies.  The tolerance around the
prescribed densities tends to zero after taking $n\to\infty$;
see~\cite[Section~1]{DemboLubetzky2018}
for the large-deviation formulation.
